
% Options for packages loaded elsewhere
\PassOptionsToPackage{unicode}{hyperref}
\PassOptionsToPackage{hyphens}{url}
\PassOptionsToPackage{dvipsnames,svgnames,x11names}{xcolor}

\documentclass[
  11pt,
  a4paper,
  notitlepage
]{report}

\usepackage{amsmath,amssymb}
\usepackage{iftex}

\ifPDFTeX
  \usepackage[T1]{fontenc}
  \usepackage[utf8]{inputenc}
  \usepackage{textcomp}
\else
  \usepackage{unicode-math}
  \defaultfontfeatures{Scale=MatchLowercase}
  \defaultfontfeatures[\rmfamily]{Ligatures=TeX,Scale=1}
\fi

\usepackage{lmodern}

\IfFileExists{upquote.sty}{\usepackage{upquote}}{}

\IfFileExists{microtype.sty}{
  \usepackage[]{microtype}
  \UseMicrotypeSet[protrusion]{basicmath}
}{}

\makeatletter
\@ifundefined{KOMAClassName}{
  \IfFileExists{parskip.sty}{
    \usepackage{parskip}
  }{
    \setlength{\parindent}{0pt}
    \setlength{\parskip}{6pt plus 2pt minus 1pt}
  }
}{
  \KOMAoptions{parskip=half}
}
\makeatother

\usepackage{xcolor}
\usepackage[margin=1in]{geometry}
\usepackage{longtable,booktabs,array}
\usepackage{calc}
\usepackage{etoolbox}

\makeatletter
\patchcmd\longtable{\par}
  {\if@noskipsec\mbox{}\fi\par}{}{}
\makeatother

\IfFileExists{footnotehyper.sty}{
  \usepackage{footnotehyper}
}{
  \usepackage{footnote}
}
\makesavenoteenv{longtable}

\setlength{\emergencystretch}{3em}

\providecommand{\tightlist}{
  \setlength{\itemsep}{0pt}
  \setlength{\parskip}{0pt}
}

% Number chapters, sections, and subsections
\setcounter{secnumdepth}{5}

% Only display chapters in the table of contents
\setcounter{tocdepth}{1}

\ifLuaTeX
  \usepackage{selnolig}
\fi

\usepackage{bookmark}
\IfFileExists{xurl.sty}{\usepackage{xurl}}{}
\urlstyle{same}

\hypersetup{
  pdftitle={Modes of Convergence of Probability Measures:
    A Uniform Empirical-Process Case Study},
  pdfauthor={Rui-Yang Zhang},
  colorlinks=true,
  linkcolor={blue},
  filecolor={Maroon},
  citecolor={Blue},
  urlcolor={blue},
  pdfcreator={LaTeX via pandoc}
}

% ================================================
% TABLE OF CONTENTS WITHOUT PAGE BREAK
% ================================================

\makeatletter
\renewcommand{\tableofcontents}{%
  \section*{\contentsname}%
  \@starttoc{toc}%
}
\makeatother

% ================================================
% TITLE
% ================================================

\title{Modes of Convergence of Probability Measures:
A Uniform Empirical-Process Case Study}
\author{Rui-Yang Zhang}
\date{November 2026}

\begin{document}

\maketitle

% ================================================
% ABSTRACT
% ================================================

\begin{abstract}
Different notions of convergence of probability measures formalise
different ways in which distributions become close. We review weak,
Wasserstein-2, total variation, and the two directions of
Kullback--Leibler convergence, together with the integral probability
metric viewpoint and maximum mean discrepancy. Using i.i.d. observations
from the uniform distribution on $[0,1]$, we compare the raw empirical
measure with a regularised histogram estimator. The former converges
weakly, in Wasserstein-2, and in maximum mean discrepancy, but remains
maximally separated in total variation and has infinite KL divergence in
both directions. The histogram converges in every considered mode. We
conclude by comparing the relative strengths of these modes, paying
particular attention to the role of compact support.
\end{abstract}

% ================================================
% TABLE OF CONTENTS
% ================================================

\vspace{0.5em}
\tableofcontents


\chapter{Notions of Convergence}\label{notions-of-convergence}

Let \(P_n\) and \(P\) be probability measures on a metric space
\((\Theta,d)\), with \(P\) the intended limit. Different notions of
convergence specify which discrepancies between \(P_n\) and \(P\) must
vanish. Unless otherwise specified, convergence below refers to a
deterministic sequence of measures; in the empirical-process example,
the resulting assertions will hold almost surely.

\section{Weak Convergence}\label{weak-convergence}

We write \(P_n\Rightarrow P\) if \[
\int f\,dP_n\longrightarrow\int f\,dP
\] for every bounded continuous function \(f:\Theta\to\mathbb R\). Weak
convergence requires agreement in the expectations of continuous
observables, but not necessarily in the probability of every measurable
event.

\section{Total Variation Convergence}\label{total-variation-convergence}

The total variation distance is \[
d_{\mathrm{TV}}(P,Q)=\sup_{A\in\mathcal B(\Theta)}|P(A)-Q(A)|.
\] When \(P\) and \(Q\) have densities \(p,q\) with respect to a common
measure, \[
d_{\mathrm{TV}}(P,Q)=\frac12\int|p-q|.
\] Thus \(d_{\mathrm{TV}}(P_n,P)\to0\) requires uniform agreement across
all measurable event probabilities.

\section{Forward and Reverse KL
Convergence}\label{forward-and-reverse-kl-convergence}

For \(P\ll Q\), define \[
D_{\mathrm{KL}}(P\|Q)=\int\log\left(\frac{dP}{dQ}\right)dP,
\] and set the divergence to \(+\infty\) otherwise. We distinguish the
two directions \[
D_{\mathrm{KL}}(P_n\|P)\to0
\qquad\text{and}\qquad
D_{\mathrm{KL}}(P\|P_n)\to0.
\] We call these \emph{forward} and \emph{reverse} KL convergence,
respectively, relative to the fixed target \(P\) (these labels are not
universal). Both control relative density discrepancies, but they are
not equivalent in general.

\section{Wasserstein-2 Convergence}\label{wasserstein-2-convergence}

For probability measures with finite second moments, \[
W_2(P,Q)^2=\inf_{\gamma\in\Gamma(P,Q)}
\int_{\Theta\times\Theta}d(x,y)^2\,d\gamma(x,y),
\] where \(\Gamma(P,Q)\) is the set of couplings with marginals \(P\)
and \(Q\). Wasserstein-2 convergence, \(W_2(P_n,P)\to0\), requires that
mass can be transported between the distributions at vanishing mean
squared cost. On \(\mathbb R^d\), it is equivalent to weak convergence
together with convergence of second moments.

\section{Integral Probability Metrics and
MMD}\label{integral-probability-metrics-and-mmd}

An integral probability metric (IPM) associated with a function class
\(\mathcal F\) is \[
d_{\mathcal F}(P,Q)=\sup_{f\in\mathcal F}
\left|\int f\,dP-\int f\,dQ\right|.
\] Total variation is an IPM with \(\mathcal F=\{f:0\le f\le1\}\), while
Wasserstein-1 corresponds to 1-Lipschitz functions (under the usual
moment assumptions). The \textbf{maximum mean discrepancy} (MMD) uses
the unit ball of a reproducing kernel Hilbert space \(\mathcal H_k\): \[
\operatorname{MMD}_k(P,Q)
=\sup_{\|f\|_{\mathcal H_k}\le1}|P f-Q f|.
\] In our case study we use the Gaussian kernel
\(k_\sigma(x,y)=\exp\{-(x-y)^2/(2\sigma^2)\}\) with fixed \(\sigma>0\).
Its MMD metrises weak convergence on \([0,1]\). KL and Wasserstein-2 are
not, in general, IPMs.

\chapter{A Uniform Empirical-Process Case
Study}\label{a-uniform-empirical-process-case-study}

Suppose \[
X_1,X_2,\ldots\overset{\mathrm{iid}}{\sim}P=\operatorname{Unif}[0,1],
\qquad p(x)=1\quad (x\in[0,1]).
\] The \textbf{raw empirical measure} and empirical cumulative
distribution function (CDF) are \[
\widehat P_n=\frac1n\sum_{i=1}^n\delta_{X_i},
\qquad
F_n(x)=\frac1n\sum_{i=1}^n\mathbf1\{X_i\le x\}.
\] The target CDF is \(F(x)=x\) for \(x\in[0,1]\). We first analyse
\(\widehat P_n\) and then introduce a smoothed counterpart
\(\widetilde P_n\). This distinction is essential: the raw empirical
distribution is discrete even when the target is continuous.

\section{A Smoothed Empirical Measure and a Central
Estimate}\label{a-smoothed-empirical-measure-and-a-central-estimate}

Let \(J_n=\lfloor n^{1/4}\rfloor\) and partition \([0,1]\) into \(J_n\)
equal bins \(B_{1,n},\ldots,B_{J_n,n}\), with endpoint conventions
chosen to form a partition. Define the bin counts \[
N_{j,n}=\sum_{i=1}^n\mathbf1\{X_i\in B_{j,n}\}.
\] We add one pseudocount per bin and define a \textbf{regularised
histogram density} by \[
\boxed{
\widetilde p_n(x)=\frac{J_n(N_{j,n}+1)}{n+J_n},
\qquad x\in B_{j,n}.
}
\] The corresponding measure is
\(\widetilde P_n(dx)=\widetilde p_n(x)\,dx\). It is normalised and
strictly positive. If \(H_n\) denotes the histogram obtained by
spreading each empirical atom uniformly over its bin, then \[
\widetilde P_n=\frac{n}{n+J_n}H_n
+\frac{J_n}{n+J_n}P.
\] The second term is a vanishing uniform component, which prevents
zero-density bins.

Our central estimate is \[
\boxed{\delta_n:=\|\widetilde p_n-1\|_\infty
=O_{\mathrm{a.s.}}\!\left(n^{-1/4}\sqrt{\log n}\right)\longrightarrow0.}
\] To verify it, write \(\Delta_n=\|F_n-F\|_\infty\). The
Dvoretzky--Kiefer--Wolfowitz (DKW) inequality and Borel--Cantelli give
\[
\Delta_n=O_{\mathrm{a.s.}}\!\left(\sqrt{\frac{\log n}{n}}\right).
\] For every bin, \(|N_{j,n}/n-1/J_n|\le2\Delta_n\). Since the uniform
population histogram has density exactly one, \[
\begin{aligned}
\delta_n
&=\max_{1\le j\le J_n}
\left|\frac{J_nN_{j,n}-n}{n+J_n}\right|\\
&\le\frac{2nJ_n}{n+J_n}\Delta_n
\le2J_n\Delta_n.
\end{aligned}
\] Because \(J_n\asymp n^{1/4}\), this tends to zero almost surely.
Unlike a non-uniform target, the uniform distribution incurs no
histogram approximation bias; only the empirical bin fluctuations
remain.

\section{Weak Convergence}\label{weak-convergence-1}

For the raw empirical measure, Glivenko--Cantelli gives \[
\Delta_n=\sup_{x\in\mathbb R}|F_n(x)-F(x)|
\overset{\mathrm{a.s.}}{\longrightarrow}0.
\] Uniform CDF convergence implies weak convergence, and therefore \[
\boxed{\widehat P_n\Rightarrow P\qquad\mathrm{a.s.}}
\] Equivalently, the strong law gives
\(n^{-1}\sum_i f(X_i)\to\mathbb E_Pf(X)\) for bounded continuous \(f\);
the CDF argument supplies a common almost-sure event.

For the smoothed measure, let \(f\) be bounded and continuous, and
define its modulus of continuity by
\(\omega_f(h)=\sup_{|x-y|\le h}|f(x)-f(y)|\). Spreading an observation
uniformly within its bin changes \(f\) by at most \(\omega_f(1/J_n)\).
From the mixture representation, \[
\left|\int f\,d\widetilde P_n-\int f\,d\widehat P_n\right|
\le\omega_f(1/J_n)
+\frac{2J_n}{n+J_n}\|f\|_\infty\longrightarrow0.
\] Together with \(\widehat P_n\Rightarrow P\) this yields \[
\boxed{\widetilde P_n\Rightarrow P\qquad\mathrm{a.s.}}
\] The proof relies only on the empirical-process law of large numbers
and the fact that smoothing becomes negligible for continuous
observables.

\section{Wasserstein-2 Convergence}\label{wasserstein-2-convergence-1}

On \([0,1]\), \(|x-y|^2\le|x-y|\); hence \(W_2^2(Q,R)\le W_1(Q,R)\). In
one dimension, \[
W_1(\widehat P_n,P)=\int_0^1|F_n(x)-x|\,dx
\le\Delta_n.
\] It follows that \[
\boxed{W_2(\widehat P_n,P)\le\sqrt{\Delta_n}
\longrightarrow0\qquad\mathrm{a.s.}}
\] For \(\widetilde P_n\), we construct an explicit coupling with
\(\widehat P_n\). With probability \(n/(n+J_n)\), pair a randomly
selected observation with a uniform point in its own bin, whose distance
is at most \(1/J_n\). Otherwise, pair a randomly selected observation
with an independent uniform point in \([0,1]\), whose distance is at
most one. Consequently, \[
W_2^2(\widetilde P_n,\widehat P_n)
\le\frac1{J_n^2}+\frac{J_n}{n+J_n}.
\] The triangle inequality then gives \[
\boxed{
W_2(\widetilde P_n,P)
\le\left(\frac1{J_n^2}+\frac{J_n}{n+J_n}\right)^{1/2}
+W_2(\widehat P_n,P)\longrightarrow0
\quad\mathrm{a.s.}}
\] Here the crucial ingredient is control of the transportation
distance, rather than control of density ratios.

\section{Total Variation
Convergence}\label{total-variation-convergence-1}

Let \(A_n=\{X_1,\ldots,X_n\}\). For every finite sample,
\(\widehat P_n(A_n)=1\) whereas \(P(A_n)=0\) almost surely. Thus \[
\boxed{d_{\mathrm{TV}}(\widehat P_n,P)=1
\qquad\text{for every }n\ \mathrm{a.s.}}
\] In contrast, \(\widetilde P_n\) and \(P\) have Lebesgue densities,
and \[
\begin{aligned}
d_{\mathrm{TV}}(\widetilde P_n,P)
&=\frac12\int_0^1|\widetilde p_n(x)-1|\,dx\\
&\le\frac12\delta_n\longrightarrow0.
\end{aligned}
\] Therefore \[
\boxed{d_{\mathrm{TV}}(\widetilde P_n,P)\longrightarrow0
\qquad\mathrm{a.s.}}
\] Smoothing removes the discrete-versus-continuous singularity, and
uniform density control is more than sufficient for TV convergence.

\section{Forward KL Convergence}\label{forward-kl-convergence}

Since \(\widehat P_n\) assigns mass to points to which \(P\) assigns
none, \(\widehat P_n\not\ll P\). Consequently, \[
\boxed{D_{\mathrm{KL}}(\widehat P_n\|P)=+\infty
\qquad\text{for every }n\ \mathrm{a.s.}}
\] For the histogram, use \(D_{\mathrm{KL}}(Q\|R)\le\chi^2(Q\|R)\) and
\(p\equiv1\): \[
\begin{aligned}
D_{\mathrm{KL}}(\widetilde P_n\|P)
&\le\int_0^1\frac{(\widetilde p_n(x)-1)^2}{1}\,dx\\
&\le\delta_n^2\longrightarrow0.
\end{aligned}
\] Hence \[
\boxed{D_{\mathrm{KL}}(\widetilde P_n\|P)\longrightarrow0
\qquad\mathrm{a.s.}}
\] The fixed reference density is bounded away from zero, so an absolute
density error also controls the relative error entering forward KL.

\section{Reverse KL Convergence}\label{reverse-kl-convergence}

Because \(P\) is nonatomic while \(\widehat P_n\) is supported on
finitely many points, \(P\not\ll\widehat P_n\), giving \[
\boxed{D_{\mathrm{KL}}(P\|\widehat P_n)=+\infty
\qquad\text{for every }n\ \mathrm{a.s.}}
\] For the histogram, the pseudocounts ensure \(\widetilde p_n>0\) for
every \(n\). Moreover, \(\delta_n\to0\) implies that eventually
\(\widetilde p_n(x)\ge1/2\) uniformly. Thus \[
\begin{aligned}
D_{\mathrm{KL}}(P\|\widetilde P_n)
&\le\int_0^1\frac{(1-\widetilde p_n(x))^2}{\widetilde p_n(x)}\,dx\\
&\le2\delta_n^2\longrightarrow0.
\end{aligned}
\] We conclude that \[
\boxed{D_{\mathrm{KL}}(P\|\widetilde P_n)\longrightarrow0
\qquad\mathrm{a.s.}}
\] Unlike forward KL, this bound requires control of the \emph{estimated
density in the denominator}. The positive pseudocounts guarantee finite
reverse KL at every \(n\), and uniform consistency supplies the eventual
lower bound.

\section{Maximum Mean Discrepancy
Convergence}\label{maximum-mean-discrepancy-convergence}

Fix a Gaussian kernel \(k_\sigma\) and let
\(\varphi(x)=k_\sigma(x,\cdot)\in\mathcal H_{k_\sigma}\). The empirical
MMD is the Hilbert-space norm \[
\operatorname{MMD}_{k_\sigma}(\widehat P_n,P)
=\left\|\frac1n\sum_{i=1}^n\varphi(X_i)
-\mathbb E_P\varphi(X)\right\|_{\mathcal H_{k_\sigma}}.
\] Since \(\|\varphi(x)\|_{\mathcal H_{k_\sigma}}^2=k_\sigma(x,x)=1\),
the Hilbert-space strong law of large numbers gives \[
\boxed{\operatorname{MMD}_{k_\sigma}(\widehat P_n,P)
\longrightarrow0\qquad\mathrm{a.s.}}
\] For the smoothed measure, the reproducing property gives
\(\|f\|_\infty\le1\) whenever \(\|f\|_{\mathcal H_{k_\sigma}}\le1\).
Hence \[
\operatorname{MMD}_{k_\sigma}(\widetilde P_n,P)
\le\int_0^1|\widetilde p_n(x)-1|\,dx
\le\delta_n\longrightarrow0.
\] Thus both measures converge in Gaussian-kernel MMD, even though the
raw empirical measure fails to converge in TV or either KL direction.

\chapter{Relative Strengths of the Convergences}
\label{relative-strengths-of-the-convergences}

We now compare the modes of convergence introduced in
Chapter 1. Our main objective is to understand which notions
imply others, why these implications hold, and under what
additional assumptions stronger conclusions may be obtained.

\section{General Implications}
\label{general-implications}

Let $P_n$ and $P$ be probability measures on a Polish
metric space $(\Theta,d)$. In general, the notions of
convergence considered here are only partially ordered:
some implications hold universally, while others require
additional assumptions on the distributions or the
underlying space.

\subsection{KL Convergence Implies Total Variation Convergence}

Recall Pinsker's inequality:
\[
d_{\mathrm{TV}}(Q,P)
\leq
\sqrt{\frac{1}{2}D_{\mathrm{KL}}(Q\|P)}.
\]
Since total variation is symmetric, we may interchange
the two arguments to obtain
\[
d_{\mathrm{TV}}(Q,P)
\leq
\sqrt{\frac{1}{2}D_{\mathrm{KL}}(P\|Q)}.
\]
Consequently, convergence in either direction of KL
implies total variation convergence:
\[
\boxed{
\begin{aligned}
D_{\mathrm{KL}}(P_n\|P)\to0
&\implies d_{\mathrm{TV}}(P_n,P)\to0,\\
D_{\mathrm{KL}}(P\|P_n)\to0
&\implies d_{\mathrm{TV}}(P_n,P)\to0.
\end{aligned}
}
\]

The converse does not generally hold. Total variation
controls absolute discrepancies in event probabilities,
whereas KL additionally penalises large relative density
discrepancies. A density can be close to another in
$L^1$ while differing by an arbitrarily large factor
on a sufficiently small region.

Furthermore, the two directions of KL convergence are
incomparable: neither implies the other. We provide
explicit examples in Section~3.2.

\subsection{Total Variation Implies Weak Convergence}

Recall the variational characterisation
\[
d_{\mathrm{TV}}(P_n,P)
=
\frac12
\sup_{\|f\|_\infty\leq1}
\left|
\int f\,dP_n-\int f\,dP
\right|.
\]
For any bounded measurable function $f$, it follows that
\[
\left|
\int f\,dP_n-\int f\,dP
\right|
\leq
2\|f\|_\infty d_{\mathrm{TV}}(P_n,P).
\]
Thus, if $d_{\mathrm{TV}}(P_n,P)\to0$, then
\[
\int f\,dP_n\longrightarrow\int f\,dP
\]
for every bounded measurable $f$. In particular,
this holds for every bounded continuous function,
which is precisely the definition of weak convergence.
Therefore,
\[
\boxed{
d_{\mathrm{TV}}(P_n,P)\to0
\implies P_n\Rightarrow P.
}
\]

The converse fails because weak convergence only
tests bounded continuous functions, whereas total
variation controls all bounded measurable functions.
In particular, weak convergence cannot necessarily
distinguish between probability masses located at
two arbitrarily close but distinct points.

\subsection{Wasserstein-2 Implies Weak Convergence}

For probability measures with finite second moments,
we have
\[
W_1(P_n,P)\leq W_2(P_n,P).
\]
By the Kantorovich--Rubinstein duality,
\[
W_1(P_n,P)
=
\sup_{\operatorname{Lip}(f)\leq1}
\left|
\int f\,dP_n-\int f\,dP
\right|.
\]
Consequently, $W_2(P_n,P)\to0$ implies convergence
of expectations for every bounded Lipschitz function.
Since bounded Lipschitz functions characterise weak
convergence on Polish spaces, we obtain
\[
\boxed{
W_2(P_n,P)\to0
\implies P_n\Rightarrow P.
}
\]

However, Wasserstein-2 convergence is stronger than
weak convergence on unbounded spaces. In fact,
for probability measures on $\mathbb R^d$ with
finite second moments,
\[
\boxed{
W_2(P_n,P)\to0
\iff
\begin{cases}
P_n\Rightarrow P,\\
\displaystyle
\int\|x\|^2\,dP_n(x)
\to
\int\|x\|^2\,dP(x).
\end{cases}
}
\]
Thus, $W_2$ convergence additionally requires
convergence of second moments. Weak convergence
alone does not control the contribution of a
vanishing amount of probability mass escaping
arbitrarily far from the origin.

\subsection{Incomparability of Wasserstein and TV/KL}

An important distinction is that Wasserstein
distances depend on the geometry of the underlying
space, whereas TV and KL are insensitive to the
distance between individual points.

For example, consider
\[
P_n=\delta_{1/n},
\qquad P=\delta_0.
\]
Clearly,
\[
W_2(P_n,P)=\frac1n\longrightarrow0,
\]
but
\[
d_{\mathrm{TV}}(P_n,P)=1
\]
for every $n$. Therefore, Wasserstein-2 convergence
does not imply TV convergence, even on compact spaces.

The converse also fails on unbounded spaces.
Indeed, even convergence in both directions of KL
need not imply Wasserstein-2 convergence.

To illustrate this, consider probability measures
on the positive integers, equipped with the usual
Euclidean distance. Let
\[
P(\{k\})=\frac{C}{k^4},
\qquad k=1,2,\ldots,
\]
where $C$ is the normalising constant, and define
\[
P_n=
\left(1-\frac1{n^2}\right)P
+\frac1{n^2}\delta_n,
\qquad n\geq2.
\]
By convexity of KL in its first argument,
\[
\begin{aligned}
D_{\mathrm{KL}}(P_n\|P)
&\leq
\frac1{n^2}
D_{\mathrm{KL}}(\delta_n\|P)\\
&=
\frac1{n^2}\log\frac1{P(\{n\})}\\
&=
O\left(\frac{\log n}{n^2}\right)
\longrightarrow0.
\end{aligned}
\]
Moreover, since
$P_n\geq(1-n^{-2})P$, we have
\[
D_{\mathrm{KL}}(P\|P_n)
\leq-\log(1-n^{-2})\longrightarrow0.
\]
Thus both KL divergences vanish.

However, writing $m_2=\mathbb E_P[X^2]$, we obtain
\[
\begin{aligned}
\mathbb E_{P_n}[X^2]
&=
\left(1-\frac1{n^2}\right)m_2
+\frac1{n^2}n^2\\
&\longrightarrow m_2+1.
\end{aligned}
\]
The second moments therefore fail to converge to
that of $P$, implying
\[
W_2(P_n,P)\not\longrightarrow0.
\]

This example demonstrates that KL convergence,
despite being stronger than TV convergence, does
not necessarily provide the moment control required
for Wasserstein convergence.

\subsection{Maximum Mean Discrepancy and Weak Convergence}

Recall that MMD is an integral probability metric
associated with the unit ball of a reproducing kernel
Hilbert space $\mathcal H_k$:
\[
\operatorname{MMD}_k(P,Q)
=
\sup_{\|f\|_{\mathcal H_k}\leq1}
|Pf-Qf|.
\]

Suppose $k(x,x)\leq K$ uniformly. By the reproducing
property,
\[
\|f\|_\infty\leq\sqrt K
\qquad
\text{whenever }\|f\|_{\mathcal H_k}\leq1.
\]
Therefore,
\[
\boxed{
\operatorname{MMD}_k(P,Q)
\leq2\sqrt K\,d_{\mathrm{TV}}(P,Q).
}
\]
Hence TV convergence implies MMD convergence
for every bounded kernel.

For the Gaussian kernel on $\mathbb R^d$,
\[
k_\sigma(x,y)
=
\exp\left\{
-\frac{\|x-y\|^2}{2\sigma^2}
\right\},
\]
MMD also metrises weak convergence:
\[
\boxed{
\operatorname{MMD}_{k_\sigma}(P_n,P)\to0
\iff P_n\Rightarrow P.
}
\]

This equivalence is specific to sufficiently rich
kernels and does not hold for arbitrary choices
of $k$. In particular, merely requiring a kernel
to distinguish probability measures does not
always ensure that its MMD metrises weak convergence.

\section{The Simplified Hierarchy on $[0,1]$}
\label{the-simplified-hierarchy-on-01}

The relationships become substantially simpler
when all probability measures are supported on a
compact metric space.

Suppose $(\Theta,d)$ has finite diameter $D$.
By the maximal coupling characterisation of
total variation, we can construct a coupling
$(X,Y)$ of $P$ and $Q$ such that
\[
\mathbb P(X\neq Y)=d_{\mathrm{TV}}(P,Q).
\]
Since $d(X,Y)\leq D$, it follows that
\[
\begin{aligned}
W_2^2(P,Q)
&\leq\mathbb E[d(X,Y)^2]\\
&\leq D^2\mathbb P(X\neq Y)\\
&=D^2d_{\mathrm{TV}}(P,Q).
\end{aligned}
\]
Consequently,
\[
\boxed{
d_{\mathrm{TV}}(P_n,P)\to0
\implies W_2(P_n,P)\to0.
}
\]

Furthermore, compact support ensures convergence
of all finite-order moments under weak convergence.
Therefore, for every finite $p\geq1$,
\[
W_p(P_n,P)\to0
\iff P_n\Rightarrow P.
\]

In particular, on $[0,1]$, the modes considered
in this note satisfy the hierarchy
\[
\boxed{
\begin{gathered}
D_{\mathrm{KL}}(P_n\|P)\to0
\quad\text{or}\quad
D_{\mathrm{KL}}(P\|P_n)\to0
\\[5pt]
\Downarrow
\\[5pt]
d_{\mathrm{TV}}(P_n,P)\to0
\\[5pt]
\Downarrow
\\[5pt]
W_2(P_n,P)\to0
\quad\Longleftrightarrow\quad
P_n\Rightarrow P
\\[5pt]
\Longleftrightarrow
\\[5pt]
\operatorname{MMD}_{k_\sigma}(P_n,P)\to0.
\end{gathered}
}
\]
Here $k_\sigma$ denotes the fixed Gaussian kernel
considered in Chapter 2.

Importantly, the two KL convergence statements
at the top of the hierarchy remain incomparable,
even on $[0,1]$. To see this, let
$P=\operatorname{Unif}[0,1]$.

First, consider the density
\[
q_n(x)=
\begin{cases}
e^{-n^2}, & x\in[0,1/n],\\[3pt]
\displaystyle
\frac{1-e^{-n^2}/n}{1-1/n},
& x\in(1/n,1].
\end{cases}
\]
Writing $Q_n(dx)=q_n(x)\,dx$, one can verify that
\[
D_{\mathrm{KL}}(Q_n\|P)\to0,
\]
whereas
\[
D_{\mathrm{KL}}(P\|Q_n)\to\infty.
\]
The reverse KL diverges because $q_n$ becomes
exponentially small on a shrinking interval.

Conversely, let $\eta_n=e^{-n^2}$ and consider
\[
r_n(x)
=
1-\frac1n+
\frac{1}{n\eta_n}
\mathbf1_{[0,\eta_n]}(x).
\]
This is also a probability density on $[0,1]$.
Writing $R_n(dx)=r_n(x)\,dx$, we have
\[
D_{\mathrm{KL}}(P\|R_n)\to0,
\]
but
\[
D_{\mathrm{KL}}(R_n\|P)\to\infty.
\]
Here the forward KL diverges because $r_n$
develops an increasingly concentrated density spike.

Both sequences nevertheless converge to $P$ in
total variation. Thus TV convergence is strictly
weaker than either KL convergence, and neither
direction of KL dominates the other.

Finally, the implication from TV to Wasserstein-2
is also strict, as demonstrated by the raw empirical
measure in Chapter 2: it converges in $W_2$ while
remaining at TV distance one from the target.

\section{Conclusions from the Case Study}
\label{conclusions-from-the-case-study}

The empirical-process example illustrates how the
abstract hierarchy translates into concrete
statistical convergence results.

\begin{center}
\small
\renewcommand{\arraystretch}{1.2}
\begin{tabular}{@{}lcc@{}}
\toprule
Mode
& Raw empirical measure
& Smoothed histogram\\
\midrule
Weak
& Converges & Converges\\
Wasserstein-2
& Converges & Converges\\
Total variation
& Distance always $1$ & Converges\\
Forward KL
& Always $+\infty$ & Converges\\
Reverse KL
& Always $+\infty$ & Converges\\
Gaussian MMD
& Converges & Converges\\
\bottomrule
\end{tabular}
\end{center}

All convergence statements in the table hold
almost surely.

The fundamental distinction between the raw
empirical measure and the smoothed histogram is
their measure-theoretic structure. Although the
empirical measure approximates the target in the
weak and Wasserstein senses, its discrete support
prevents convergence in TV or either KL direction.

Smoothing resolves this singularity. Moreover,
the uniform consistency of the histogram density,
together with its eventual positive lower bound,
provides sufficient control to establish
convergence in both KL directions.

We conclude with two observations. First,
the strength of a convergence result depends
on the class of discrepancies it controls:
TV concerns event probabilities, KL concerns
relative densities, and Wasserstein distances
additionally reflect the geometry of probability
mass. Second, the hierarchy depends on the
properties of the underlying space. Compact
support makes weak and Wasserstein convergence
equivalent, while on unbounded spaces additional
moment or transport assumptions are generally
required to relate them.

\section{References}\label{references}

\begin{itemize}
\tightlist
\item
  Billingsley, P. (1999). \emph{Convergence of Probability Measures},
  2nd ed.~Wiley.
\item
  van der Vaart, A. W. (1998). \emph{Asymptotic Statistics}. Cambridge
  University Press.
\item
  Villani, C. (2009). \emph{Optimal Transport: Old and New}. Springer.
\item
  Gretton, A., Borgwardt, K. M., Rasch, M. J., Schoelkopf, B., and
  Smola, A. (2012). A kernel two-sample test. \emph{Journal of Machine
  Learning Research}, 13, 723--773.
\end{itemize}

\end{document}
